\documentclass[12pt]{article}

\usepackage[T2A]{fontenc}
\usepackage[utf8]{inputenc}
\usepackage[russian]{babel}

\usepackage{amsmath,amssymb,mathtools}
\usepackage{geometry}
\geometry{margin=2.5cm}

\setlength{\parindent}{0pt}
\setlength{\parskip}{4pt}

\newcommand{\dd}{\,\mathrm{d}}
\newcommand{\pd}[2]{\frac{\partial #1}{\partial #2}}
\newcommand{\pdd}[2]{\frac{\partial^2 #1}{\partial #2^2}}

\title{Вывод уравнения Блэка--Шоулза}
\author{Ширяев Е.А}
\date{}

\begin{document}
\maketitle

\section{Портфель без покрытия}

Пусть цена базового актива $S_t$ является случайным процессом,
а стоимость производного инструмента обозначим через $V(S,t)$.

Рассмотрим портфель, состоящий из длинной позиции в производном инструменте
и короткой позиции в базовом активе. На данном этапе будем считать,
что портфель сформирован без покрытия.

Стоимость портфеля:
\begin{equation}
\Pi_t = V(S_t,t) - S_t.
\end{equation}

Изменение стоимости портфеля:
\begin{equation}
\dd \Pi_t = \dd V(S_t,t) - \dd S_t.
\end{equation}

\section{Динамика базового актива}

Предположим, что цена акции описывается геометрическим броуновским движением:
\begin{equation}
\dd S_t = \mu S_t \dd t + \sigma S_t \dd W_t,
\end{equation}
где $\mu$ --- дрейф, $\sigma$ --- волатильность,
$W_t$ --- винеровский процесс.

Таким образом случайная составляющая вносит неопределенность в наш портфель через
короткую позицию по базовому активу и длинную позицию по производному финансовому
инструменту.

\section{Динамика производного инструмента}

Лемма Ито позволяет нам найти дифференциал функции, зависящей от времени
и стохастического процесса:
\begin{equation}
\dd V = \pd{V}{t}\dd t + \pd{V}{S}\dd S + \frac{1}{2}\pdd{V}{S}(\dd S)^2.
\end{equation}

Мы считаем, что небольшое изменение стоимость нашего дериватива на акции определяется выражением (4).

Найдём $(\dd S)^2$ из динамики изменения цены:
\[
(\dd S)^2 = (\mu S \dd t + \sigma S \dd W)^2.
\]

Раскрываем:
\[
(\dd S)^2 =
\mu^2 S^2 \dd t^2
+ 2\mu\sigma S^2 \dd t\,\dd W
+ \sigma^2 S^2 (\dd W)^2.
\]

Используя правила стохастического исчисления:
\[
(\dd W)^2=\dd t,\qquad
\dd W\,\dd t=0,\qquad
\dd t^2=0,
\]
получаем:
\begin{equation}
(\dd S)^2=\sigma^2 S^2 \dd t.
\end{equation}

Подставляя:
\begin{equation}
\dd V =
\pd{V}{t}\dd t
+ \pd{V}{S}(\mu S \dd t + \sigma S \dd W)
+ \frac{1}{2}\sigma^2 S^2 \pdd{V}{S}\dd t.
\end{equation}

\section{Источник случайности}

На первый взгляд может показаться, что две составляющие уравнения вносят
элемент случайности в изменение стоимости дериватива:
\[
\frac{1}{2}\frac{\partial^2 V}{\partial S^2}(\dd S)^2,
\qquad
\frac{\partial V}{\partial S}\dd S.
\]

Рассмотрим первый член. Используя соотношение
\[
(\dd S)^2=\sigma^2 S^2 \dd t,
\]
которое следует из формулы Ито (см. (4)–(5)), получаем, что данный член имеет
порядок $\dd t$ и, следовательно, является детерминированным.

Второй член содержит приращение процесса $S_t$. Подставляя динамику
базового актива
\[
\dd S = \mu S \dd t + \sigma S \dd W,
\]
видим, что именно здесь возникает стохастическая составляющая,
пропорциональная $\dd W$.

Таким образом, единственным источником случайности в динамике стоимости
дериватива является член
\[
\sigma S \frac{\partial V}{\partial S}\dd W.
\]

\section{Введение покрытия}

Поскольку портфель без покрытия остаётся случайным,
возникает возможность выбрать количество акций таким образом,
чтобы устранить стохастическую составляющую.

Введём количество акций $\Delta_t$ и рассмотрим портфель:
\[
\Pi_t = V(S_t,t) - \Delta_t S_t.
\]

Тогда:
\[
\dd \Pi_t = \dd V - \Delta_t \dd S_t.
\]

Подставляя выражения:
\[
\dd \Pi =
\Bigl(\pd{V}{t} + \mu S \pd{V}{S} + \frac{1}{2}\sigma^2 S^2 \pdd{V}{S}\Bigr)\dd t
+ \sigma S \pd{V}{S}\dd W
- \Delta_t(\mu S\dd t + \sigma S\dd W).
\]

Сгруппируем случайные члены:
\[
\dd \Pi =
\Bigl(\pd{V}{t} + \mu S \pd{V}{S}
+ \frac{1}{2}\sigma^2 S^2 \pdd{V}{S}
- \Delta_t\mu S\Bigr)\dd t
+ \sigma S(\pd{V}{S}-\Delta_t)\dd W.
\]

Случайная составляющая исчезает при выборе
\begin{equation}
\Delta_t=\pd{V}{S}.
\end{equation}

Таким образом удаётся полностью устранить стохастическую риск.
Портфель становится локально безрисковым:
\begin{equation}
\dd \Pi =
\Bigl(\pd{V}{t} + \frac{1}{2}\sigma^2 S^2 \pdd{V}{S}\Bigr)\dd t.
\end{equation}

Поскольку портфель $\Pi_t$ не содержит стохастической составляющей,
его динамика становится детерминированной. В условиях отсутствия
арбитража любой безрисковый портфель должен приносить безрисковую
процентную ставку $r$. В противном случае возникла бы возможность
получать гарантированную прибыль без риска.

Следовательно, динамика такого портфеля должна удовлетворять
уравнению роста безрискового актива.

\section{Уравнение Блэка--Шоулза}

Безрисковый портфель должен приносить безрисковую доходность $r$:
\begin{equation}
\dd \Pi = r\Pi \dd t.
\end{equation}

Заметим, что параметр $\mu$, характеризующий ожидаемую доходность акции,
исчезает из дальнейших вычислений. Это означает, что цена дериватива
не зависит от реальной ожидаемой доходности актива. Эквивалентно можно
сказать, что вычисления проводятся в риск-нейтральной мере $Q$, в которой
динамика базового актива имеет вид
\[
\dd S = r S \dd t + \sigma S \dd W^Q .
\]

Подставляя портфель:
\[
\dd \Pi = r\bigl(V - S\pd{V}{S}\bigr)\dd t.
\]

Приравнивая:
\[
\pd{V}{t} + \frac{1}{2}\sigma^2 S^2 \pdd{V}{S}
= r\bigl(V - S\pd{V}{S}\bigr),
\]
получаем:
\begin{equation}
\boxed{
\pd{V}{t}
+ \frac{1}{2}\sigma^2 S^2 \pdd{V}{S}
+ rS\pd{V}{S}
- rV = 0
}
\end{equation}

Это уравнение Блэка--Шоулза.

\section{Закрытая форма решения}

Итак, мы получили уравнение Блэка--Шоулза в дифференциальной форме.
Что делать дальше?

Часто мы сталкиваемся с таким уравнением:
\[
C = S\cdot N(d_1) - K e^{-rT} N(d_2),
\]
\[
d_1 = \frac{\ln\left(\frac{S}{K}\right) + \left(r + \frac12\sigma^2\right)T}{\sigma\sqrt{T}},
\qquad
d_2 = d_1 - \sigma\sqrt{T}.
\]

Это уравнение Блэка--Шоулза для европейского
колл опциона (без дивидендов).

Что здесь к чему: первый вопрос, которым стоит задаться,
глядя на это уравнение.

Мы знаем, что для исполнения опционного контракта
должно выполняться условие $(S_T - K)^+$, другими словами
\begin{equation}
S_T - K > 0 \Rightarrow S_T > K.
\end{equation}

На этом моменте стоит задуматься и прологарифмировать выражение.
Получим
\[
\ln(S_T) - \ln(K).
\]

Теперь нужно вернуться назад и вспомнить,
что мы полагали, что ценообразование исходит
из геометрического броуновского движения.

В реальной (физической) мере $P$ динамика актива имеет вид
\[
\dd S = \mu S \dd t + \sigma S \dd W.
\]

Однако при оценке деривативов мы работаем
в риск-нейтральной мере $Q$, в которой ожидаемая
доходность актива равна безрисковой ставке $r$.
В этой мере динамика принимает вид
\[
\dd S = r S \dd t + \sigma S \dd W^Q.
\]

Интегрируя это уравнение, получаем
\[
\ln(S_T) = \ln(S_0) + \left(r - \frac12\sigma^2\right)T + \sigma W_T.
\]

Получим:
\[
\ln(S_0) + \left(r - \frac12\sigma^2\right)T + \sigma W_T > \ln(K),
\]
тогда
\[
\ln\left(\frac{K}{S_0}\right)
<
\left(r - \frac12\sigma^2\right)T + \sigma W_T.
\]

Так как $W_T\sim N(0,T)$, нормируем:
\[
W_T=\sqrt{T}\,Z, \qquad Z\sim N(0,1).
\]

Подставим:
\[
\ln\left(\frac{K}{S_0}\right)
<
\left(r - \frac12\sigma^2\right)T + \sigma\sqrt{T}Z.
\]

Тогда:
\[
\sigma\sqrt{T}Z >
\ln\left(\frac{K}{S_0}\right)
-
\left(r - \frac12\sigma^2\right)T,
\]
\[
Z >
\frac{1}{\sigma\sqrt{T}}
\left(
\ln\left(\frac{K}{S_0}\right)
-
\left(r - \frac12\sigma^2\right)T
\right).
\]

На этом моменте остановимся подробнее.
Ранее мы говорили, что $Z\sim N(0,1)$.

Для понимания необходимо ответить на вопрос:
что такое $N(x)$? $N(x)=P(Z\le x)$.

Подразумевается, что это площадь под кривой нормального распределения
слева от точки $x$. В нашем случае нас интересует,
чтобы call option был исполнен, поэтому нас интересует площадь,
находящаяся справа от точки.

Обозначим эту точку как $d_2$, тогда
\[
P(Z>-d_2).
\]

Перевернём выражение:
\[
Z >
\frac{1}{\sigma\sqrt{T}}
\left(
\ln\left(\frac{S_0}{K}\right)
+
\left(r - \frac12\sigma^2\right)T
\right)
= -d_2.
\]

Другими словами:
\begin{equation}
P(S_T>K)=N(d_2).
\end{equation}

Мы помним, что находимся в риск-нейтральной мере,
тогда цена опциона:
\begin{equation}
C_0=e^{-rT}\mathbb{E}^Q[(S_T-K)^+].
\end{equation}

Разберём глубже:
\begin{equation}
C_0=e^{-rT}\mathbb{E}^Q[S_T\mathbf{1}_{S_T>K}-K\mathbf{1}_{S_T>K}],
\end{equation}

\begin{equation}
C_0=e^{-rT}
\left(
\mathbb{E}^Q[S_T\mathbf{1}_{S_T>K}]
-
K P(S_T>K)
\right).
\end{equation}

Вторую часть мы уже ввели и знаем,
как с ней работать.

\[
\mathbb{E}^Q[S_T\mathbf{1}_{S_T>K}],\qquad
S_T \text{ мы уже знаем как работает:}
\]
\[
S_T=S_0 e^{(r-\frac12\sigma^2)T+\sigma\sqrt{T}Z}.
\]

Это сразу даёт ожидание:
\[
S_0 e^{(r-\frac12\sigma^2)T}
\mathbb{E}^Q\left[e^{\sigma\sqrt{T}Z}\mathbf{1}_{Z>-d_2}\right].
\]

Отсюда следует интеграл:
\[
\int_{-d_2}^{\infty}
e^{\sigma\sqrt{T}z}
\frac{1}{\sqrt{2\pi}}e^{-z^2/2}\dd z.
\]

Для удобства обозначим:
\[
\theta=\sigma\sqrt{T}.
\]

\[
\int_{-d_2}^{\infty}
\frac{1}{\sqrt{2\pi}}
e^{\theta z-\frac12 z^2}\dd z.
\]

Выполним выделение полного квадрата:
\[
e^{\frac12\theta^2}
\int_{-d_2-\theta}^{\infty}
\frac{1}{\sqrt{2\pi}}e^{-q^2/2}\dd q.
\]

Замена:
\[
q=z-\theta.
\]

Заметим, что
\[
d_1=d_2+\sigma\sqrt{T}=d_2+\theta.
\]

Поэтому нижний предел интегрирования преобразуется:
\[
-d_2-\theta=-d_1.
\]

Тогда
\[
\int_{-d_2-\theta}^{\infty}\phi(q)\dd q
=
P(Z>-d_1)=N(d_1).
\]

В итоге:
\[
\mathbb{E}^Q[S_T\mathbf{1}_{S_T>K}]
=
S_0 e^{rT}N(d_1).
\]

Вернёмся:
\begin{equation}
\fbox{$\displaystyle
\begin{aligned}
C_0 &= S_0N(d_1) - K e^{-rT}N(d_2),\\
\text{где}\quad d_1 &= \frac{\ln\left(\frac{S_0}{K}\right) + \left(r+\frac12\sigma^2\right)T}{\sigma\sqrt{T}},\\
d_2 &= d_1 - \sigma\sqrt{T}.
\end{aligned}
$}
\end{equation}
\end{document}
